\chapter*{摘要}
\addcontentsline{toc}{chapter}{摘要}

工业现场的日志诊断本质上不是单一分类问题，而是在连续事件流、有限边缘算力、不稳定网络和受限云端预算之间分配分析能力的问题。纯规则方案难以应对模板漂移与未知故障，全量上云又会把现场响应绑定于网络、费用和隐私边界。我们设计并实现了一套面向工业多设备日志的事件级云边协同系统CoTriage，其核心设计思想是不预先固定“边缘”或“云端”为唯一执行位置，而是依据每条日志事件的风险、认知不确定性与实时系统状态，动态决定分析深度和处理时机。
该系统首先构建统一事件语义，并将规则、可审计模板知识、MiniLM ONNX 轻量分类器和边缘小语言模型组织为异构分层推理链。确定性异常在低成本路径提前收敛，开放性异常逐级获得更强语义能力，各层在同一输出契约中保留证据、置信度、模型与规则版本，使速度、能力和可解释性不再相互割裂。规则、模板与模型配置可热更新并原子切换，失败时保留上一可用快照，从而将模型推理能力转化为可治理的基础设施。

CoTriage的核心创新是风险感知、不确定性校准、策略意图约束的自适应云边路由机制 RAISE。该机制将“事件是否值得上云”与“当前是否具备发送条件”分离建模。前者联合严重度、校准不确定性、模板新颖度、历史边云分歧、跨设备影响与推理路径风险评估云端增益，后者综合网络健康、队列压力、调用预算、限流和隐私策略形成动态门槛，并输出 \texttt{LOCAL}、\texttt{CLOUD}、\texttt{DEFER} 三类可解释决策。由此，高频确定性事件可在边缘闭环，高价值事件获得按需深诊；当云链路受限时，本地判断持续运行，延期任务则由持久化 Outbox 按风险优先级重放，避免弱网将“暂时不可发送”演化为“永久失去诊断机会”。

云端并非简单复判单条日志，而是在 Token 预算内聚合时间邻近、同源、同模板和跨设备关联事件，经脱敏后生成根因候选、处置建议与结构化解释。诊断结果通过幂等边界持久化，并被转换为纠错、可执行信息增量和高风险确认等反馈信号，更新模板统计与在线增益状态；统一事件标识、状态快照和全链路观测则使每次推理、路由、延期和诊断均可追溯。

CoTriage 的价值在于建立一套可度量、可约束、可演进的大模型云边协同推理机制：把云端能力集中用于少量真正困难或高风险的事件，把实时性与基本业务连续性留在现场，并以稳定契约支持日志来源、模型 Provider、节点规模和应用场景的扩展。该方案为工业生产服务、数据中心及分布式基础设施提供了兼顾实时响应、诊断深度、资源效率与弱网韧性的智能运维路径。

\noindent\textbf{关键词：}云边协同；工业日志；分层推理；RAISE 动态路由；不确定性校准；弱网韧性；根因诊断
\clearpage
